\printedpage{189}
% head: A. Grothendieck : Dérivateurs, Ch. XVIII, Édition des Universités de Montpellier II et Paris VII

\begin{enumerate}
\item[\textit{b$'$)}] \textit{$\underline{T}$ is $\underline{L}$-stable in $\cM$.}
\end{enumerate}

\noindent\textbf{Corollary 4.16.7.10.} \textit{Let $\cM$ be an accessible category. The following conditions are equivalent} (\footnote[202]{Example where these conditions are satisfied: $\cM = (k\text{-comm. alg.})$, $k$ an infinite commutative field, $\mathfrak{U}$ being the smallest universe such that $\mathbf{N}\in\mathfrak{U}$.}).
\begin{enumerate}
\item[\textit{a)}] \textit{There exists a closed subcategory of $\cM$, not accessible in $\cM$.}
\item[\textit{b)}] \textit{There exists an open subcategory in $\cM$, not accessible in $\cM$.}
\item[\textit{c)}] \textit{There exists a transfinite sequence}
\[
\underline{x} \ = \ (x_\alpha)_{\alpha\in I}
\]
\textit{of objects of $\cM$, indexed by the set $I$ of ordinals of the reference universe $\mathfrak{U}$, such that we have}
\[
\Hom(x_\alpha,x_\beta) \ = \ \emptyset \qquad \text{\textit{for} } \alpha<\beta .
\]
\end{enumerate}

\noindent Indeed, the equivalence of a) and b) is clear from the preceding corollary, and that of a) and c) is none other than 4.16.7.4.

\noindent\textsc{Proof of} 4.16.7.8. We shall have to use the

\noindent\textbf{Lemma 4.16.7.11.} \textit{Let $\cM$ be an accessible category, $\pi$ a useful cardinal for $\cM$, $\underline{S}$ a full subcategory of $\cM$, $\underline{S}(\pi)\defeq\underline{S}\cap\cM_\pi$, whence a fully faithful functor}
\[
\varphi \ : \ \Ind_\pi(\underline{S}(\pi)) \longrightarrow \cM \ .
\]

\msp{page 489}
\begin{enumerate}
\item[\textit{a)}] \textit{Let $x\in\Ob\cM$. For $x$ to belong to the essential image of the functor $\psi$, it is necessary and sufficient that the full subcategory $\underline{S}(\pi)/x$ of the category $\cM_\pi/x$ (which is filtering and even large relative to $\pi$) be cofinal (hence also large relative to $\pi$).} % [sic] \psi -> \varphi, the functor of the lemma is called \varphi in the display above
\item[\textit{b)}] \textit{For $\underline{S}$ to be $\pi$-accessible in $\cM$, it is necessary and sufficient that for every $x\in\Ob\cM$, we have}
\[
x \ \in \ \Ob\underline{S} \qquad\iff\qquad \underline{S}(\pi)/x \text{ \textit{is cofinal in} } \cM_\pi/x
\]
\textit{(or again, that $\underline{S}$ satisfy $L_\pi$, and that for every $x$ in $\cM$, we have}
\[
x \ \in\Ob\underline{S} \qquad\Longrightarrow\qquad \underline{S}(\pi)/x \text{ \textit{cofinal in} } \cM_\pi/x\text{\textit{.)}}
\]
\end{enumerate}

\noindent\textsc{Proof}. Part b) is a tautological consequence of a), which I shall prove. To say that $x$ is in the essential image of $\underline{S}$, means that
\[
x \ \simeq \ \varinjlim_{i\in I} x_i \ , \tag*{$(*)$}
\]
with $I$ large relative to $\pi$ and the $x_i$ in $\underline{S}(\pi)$. We know that $(*)$ implies ($I$ being filtering, and $\cM$ being equivalent to $\Ind_\pi(\cM)$ since $\pi$ useful for $\cM$) that $I\longrightarrow\cM_\pi/x$ is cofinal (4.4.2, page 79). Now this functor factors through the fully faithful inclusion functor $\underline{S}(\pi)/x\incl\cM_\pi/x$, which is cofinal if and only if it satisfies condition F 1. It is immediate that this condition is satisfied when it is satisfied for the composite functor in question, which is indeed the case
\printedpage{190}
% head: A. Grothendieck : Dérivateurs, Ch. XVIII, Édition des Universités de Montpellier II et Paris VII
\msp{page 490}
by hypothesis.

\noindent Conversely, suppose that $\underline{S}(\pi)/x$ is cofinal in $\cM_\pi/x$. Since we have
\[
x \xleftarrow{\ \sim\ } \varinjlim_{\cM_\pi/x} x_i \ ,
\]
$\pi$ being useful for $\cM$, we will also have by cofinality
\[
x \xleftarrow{\ \sim\ } \varinjlim_{\underline{S}(\pi)/\pi} x_i \ , % [sic] \underline{S}(\pi)/\pi -> \underline{S}(\pi)/x, the cofinal subcategory is \underline{S}(\pi)/x
\]
now $\cM_\pi$ being large relative to $\pi$, the same holds for the cofinal full subcategory $\underline{S}(\pi)/x$. Thus $x$ is in the essential image of $\Ind_\pi\underline{S}(\pi)$, q.e.d.

\noindent\textbf{Lemma 4.16.7.12.} \textit{Let $I$ be a filtering category, $J$ a closed subcategory, $K$ the complementary open subcategory. We then have the equivalences}
\[
J \ \neq \ \emptyset \iff J \text{ \textit{cofinal} } \iff K \text{ \textit{not cofinal} } \iff K \ \neq \ I \ .
\]

\noindent\textsc{Proof} immediate.

\noindent I return to the proof of 4.16.7.8. We apply Lemma 4.16.7.11 to $\cM$, $\pi$, and to the two
\msp{page 491}
subcategories $\underline{S}$, $\underline{T}$ of $\cM$, noting that for $x\in\Ob\cM$, $\underline{S}(\pi)/x$ is a closed subcategory of $\underline{S}(\pi)$, whose open complement is $\underline{T}(\pi)/x$, so that we can apply Lemma 4.16.7.12. We find:
\[
\begin{aligned}
\underline{S}\ \pi\text{-accessible} \quad &\overset{4.16.7.11}{\iff} \quad \forall\ x\ \in\ \Ob\cM,\ \text{we have } (x\ \in\ \Ob\underline{S}\iff\underline{S}(\pi)/x \text{ cofinal in}\\
&\qquad\quad \cM_\pi/x)\\
&\iff \quad \forall\, x\in\Ob\cM,\ \text{we have } \Bigl( \qquad \underbrace{x\notin\Ob\underline{S}}_{\text{(i.e. $x\in\Ob\underline{T}$, by 4.16.7.12)}}\\
&\qquad\qquad\qquad \iff \qquad \underbrace{\underline{S}(\pi)/x \text{ not cofinal in } \cM_\pi/x}_{\text{(i.e. $\underline{T}(\pi)/x$ cofinal in $\cM_\pi$, by 4.16.7.12)}} \ \Bigr)\\
&\overset{4.16.7.11}{\iff} \quad \underline{T} \text{ is } \pi\text{-accessible.}
\end{aligned}
\]
We thus have the equivalence of a) and b). On the other hand, $\underline{S}(\pi)/x$ cofinal in $\cM_\pi/x$ means simply $\underline{S}(\pi)/x\neq\emptyset$ (4.16.7.12), so $\underline{S}$ $\pi$-accessible means that $\underline{S}$ consists of the $x\in\Ob\cM$ such that $\underline{S}(\pi)\neq\emptyset$, which also means that $\underline{S}$ is the closed subcategory in $\cM$ generated by $\underline{S}(\pi)$, whence a $\iff$ a$'$. On the other hand, the equivalence b $\iff$ b$'$ is none other than 4.6.7.6 (last formulation). % [sic] 4.6.7.6 -> 4.16.7.6, the surrounding references are all to 4.16.7.x

\msp{page 492}

\section*{4.17\quad Return to paratoposes, and to the operations $\cM\boxtimes\cN$, $\underline{\operatorname{Hom}}_{!}(\cM,\cN)$}

\noindent\texttt{[This paragraph does not exist.]}

\printedpage{191}
% head: A. Grothendieck : Dérivateurs, Ch. XVIII, Édition des Universités de Montpellier II et Paris VII

\msp{page 1}

\section*{Appendix}

\subsection*{Ultraproducts of fields, and accessible categories}

\noindent Let $(k_\xi)_{\xi\in E}$ be a family of fields, indexed by a set $E$. $\Bigl($ Let $\Phi$ be a nontrivial filter in $E$, i.e. a subset $\Phi\subset\mathfrak{P}(E)$ such that:
\begin{enumerate}
\item[a)] $\emptyset\notin\Phi$, i.e. $A\in\Phi\Longrightarrow A\neq\emptyset$.
\item[b)] $A,B\in\Phi\Longrightarrow A\cap B\in\Phi$.
\item[c)] If $A\in\Phi$, $A\subset B$, then $B\in\Phi$.
\item[d)] $\Phi\neq\emptyset$ ($\iff E\in\Phi$).
\end{enumerate}
(This amounts to saying that $\Phi$ is a closed subset of the ordered set $\mathfrak{P}(E)$, \textit{left filtering}, and does not contain the least element of $\mathfrak{P}(E)$, i.e. not equal to all of $\mathfrak{P}(E)$.) $\Bigr)$ For every subset $J$ of $E$, put
\[
k_J \ = \ \prod_{\xi\in J} k_\xi \ ,
\]
we thus have an inverse system
\[
(k_J)_{J\in\mathfrak{P}(E)} \ , \qquad \text{i.e. } J\subset J' \text{ gives } k_{J'}\longrightarrow k_J \ ,
\]
and the system
\[
(k_J)_{J\in\Phi}
\]
is a direct system of rings, indexed by the filtering ordered set $\Phi^o$. Put % [sic] d'annaux -> d'anneaux, misspelling
\[
k_\Phi \ = \ \varinjlim_{J\in\Phi} k_J \ .
\]
If $\Psi$ is a filter finer than $\Phi$, i.e.
\msp{page 2}
such that $\Psi\supset\Phi$, we find a canonical homomorphism
\[
k_\Phi \longrightarrow k_\Psi \ ,
\]
so if $\cF(E)=\cF$ is the set of filters on $E$, ordered by inclusion, we find a direct system
\[
\Phi \longmapsto k_\Phi \ : \ \cF \longrightarrow \Ann \ .
\]
We have not yet used that the $k_\xi$ are fields, in fact it would suffice that they be objects of a category $\cA$ stable under small products and under small filtering $\varinjlim$. Take now $\cA=\Ann$ (category of unital rings, not necessarily commutative), and suppose that the $k_\xi$ are fields. % [sic] produis -> produits, misspelling

\printedpage{192}
% head: A. Grothendieck : Dérivateurs, Ch. XVIII, Édition des Universités de Montpellier II et Paris VII

\noindent\textbf{Proposition 1.} \textit{Suppose that the $k_\xi$ ($\xi\in E$) are} fields (\footnote[203]{NB This condition is necessary if one wants $k_\Phi$ to be a field for every ultrafilter.}). \textit{For $k_\Phi$ to be a field, it is necessary and sufficient that $\Phi$ be an ultrafilter, i.e. that for every decomposition $E=A\amalg B$, we have $A\in\Phi$ or $B\in\Phi$.} (NB If $A\in\Phi$, we have $B\notin\Phi$, since otherwise, we would have $A\cap B\in\Phi$, hence $\emptyset\in\Phi$, absurd.)

\noindent Necessity: Let, for every subset $A$ of $E$,
\[
f_A\in k_E \qquad \text{be defined by} \qquad f_A(\xi) \ = \ \left\{
\begin{array}{lll}
1_{k_\xi} & \text{if} & \xi\in A\\[1ex]
0_{k_\xi} & \text{if} & \xi\in E\backslash A \ .
\end{array}\right.
\]

\msp{page 3}
\noindent Let
\[
f_{A,\Phi} \ \in k_\Phi
\]
be the image of $f_A$ under the canonical homomorphism $k_E\longrightarrow k_\Phi$. If $E=A\amalg B$, we will have
\[
f_A \ + \ f_B \ = \ 1 \ , \qquad f_A\, f_B \ = \ 0 \ ,
\]
whence
\[
f_{A,\Phi} \ + \ f_{B,\Phi} \ = \ 1 \ , \qquad f_{A,\Phi}\, f_{B,\Phi} \ = \ 0 \ .
\]
If $k_\Phi$ is a field, this implies that we have $f_{A,\Phi}=0$ or $f_{B,\Phi}=0$. But if e.g. $f_{A,\Phi}=0$, we see that this means that there exists a $J\in\Phi$ such that $f_A|J=0$, which also means that $J\subset B$. Likewise, if $f_{B,\Phi}=0$, we will have $J\subset A$. Thus $\Phi$ is an ultrafilter.

\noindent Conversely, suppose that $\Phi$ is an ultrafilter. Let us prove that $k_\Phi$ is a field, i.e. that for $f\in k_\Phi$, $f\neq\emptyset\Longrightarrow f$ invertible. Note that $f$ comes from an element of a $k_J$ ($J\in\Phi$), hence is the canonical image of a $g\in k_E$. To say that $f=0$ means that $\exists\ J\in\Phi$ such that $g|J=0$. The opposite $f\neq0$ therefore means that for every $J\in\Phi$, we have $g|J\neq\emptyset$.

\msp{page 4}
\noindent Let then $A$ be the set of $\xi\in E$ such that $g(\xi)\neq\emptyset$, $B=E\backslash A$ the set of $\xi\in E$ such that $g(\xi)\neq\emptyset$. We thus have $B\notin\Phi$ since $f\neq0$, and since $\Phi$ is an ultrafilter, we therefore have $A\in\Phi$. But let $g'\in k_E$ be defined by $g'(\xi)=f(\xi)^{-1}$ if $\xi\in A$, $g'(\xi)=1$ if $\xi\in B$, and let $f'$ be the image of $g'$ in $k_\Phi$. We then see that we have % [sic] B = E\A ... tels que g(\xi) \neq \emptyset -> g(\xi) = 0, the complement of A is the set where g vanishes
\[
f\ f' \ = \ f' \ = \ f \ = \ 1 \ ,
\]
which indeed shows that $k_\Phi$ is a field.

\noindent Suppose now that all the $k_\xi$ ($\xi\in E$) are equal to one and the same field $k$. Then the $k_J$ form an inverse system of $k$-algebras, the $k_\Phi$ a direct system of $k$-algebras. If $k$ is commutative, the same holds for the $k_J$, $k_\Phi$, and in particular for $\Phi$ an ultrafilter, $k_\Phi$ is an \textit{extension} of the field $k$.

\noindent Let us return to the general case, but the $k_\xi$ commutative. Let $\cF_0\subset\cF$ be the set of ultrafilters on $E$. Note that for every $\Phi\in\cF$, $k_E\longrightarrow k_\Phi$ is surjective. Thus if $\Phi\in\cF$, $k_\Phi$ identifies with a field quotient of the ring $k_E$.

\printedpage{193}
% head: A. Grothendieck : Dérivateurs, Ch. XVIII, Édition des Universités de Montpellier II et Paris VII

\msp{page 5}
\noindent For every ring $\cA$, let $\Max(\cA)$ be the maximal spectrum, i.e. the set of maximal ideals of $\cA$. We thus find a canonical map
\[
\left\{
\begin{array}{lll}
\cF_0 & \longrightarrow & \Max(k_E)\\[1ex]
\ \Phi & \longmapsto & \Ker(k_E\longrightarrow k_\Phi) \ \defeq \ \mathfrak{m}_\Phi \ .
\end{array}\right. \tag*{$(*)$}
\]

\noindent\textbf{Proposition 2.} \textit{The preceding map is} bijective.

\noindent Let us first prove that it is injective, this follows at once from the

\noindent\textbf{Corollary 1} (\footnote[204]{\textbf{NB} For arbitrary $\Phi\in\cF$, $\mathfrak{m}_\Phi\defeq\Ker(\chi_\Phi : k_E\longrightarrow k_\Phi)$ determines $\Phi$, since
\[
A \ \in \ \Phi \qquad\iff\qquad \chi_\Phi(f_A) \ = \ 1
\]
(but one must assume the $K_\xi\neq0$, i.e. $0_\xi\neq1_\xi\ \ldots$).}). \textit{Let, for $\Phi\in\cF_0$,}
\[
\chi_\Phi \ : \ k_E \longrightarrow k_\Phi
\]
\textit{be the canonical map. Then for $A\in\mathfrak{P}(E)$, we have}
\[
A \ \notin \ \Phi \qquad\iff\qquad \chi_\Phi(f_A) \ = \ 0 \ , \quad \text{\textit{i.e.} } f_A\in\mathfrak{m}_\Phi
\]
\textit{(i.e. $A\in\Phi\iff\chi_\Phi(f_A)\neq0$, i.e. $f_A\notin\mathfrak{m}_\Phi$)} (\footnote[205]{\textbf{NB} We thus have
\[
(\chi_\Phi(f_A) \ \neq \ 0) \qquad\iff\qquad (\chi_\Phi(f_A) \ = \ 1) \ .
\]}).

\noindent Indeed, $\chi_\Phi(f_A)=0$ means that there exists $B\in\Phi$ such that $f_A|B=0$, i.e. such that $B\subset E\backslash A$; but this is equivalent to $E\backslash A\in\Phi$, hence to $A\notin\Phi$ since $\Phi$ is an ultrafilter, q.e.d.

\msp{page 6}
\noindent Remains to prove that the map $(*)$ is surjective. So let $\mathfrak{m}$ be a maximal ideal of $k_E$, and let
\[
\Phi \ \defeq \ \{A\in\mathfrak{P}(E)\mid\chi(f_A) \ = \ 1\} \ ,
\]
where $\chi : k_E\longrightarrow k_E/\mathfrak{m}=\kappa(\mathfrak{m})$ is the canonical map. I say that $\Phi$ is an ultrafilter. Indeed:
\begin{enumerate}
\item[a)] $\emptyset\notin\Phi$ since $f_\emptyset=0$, whence $\chi(f_\emptyset)=0\neq1$ in $k_\Phi$.
\item[b)] If $A,B\in\Phi$, we have $A\cap B\in\Phi$, since $f_{A\cap B}=f_A\cdot f_B$, whence $\chi(f_{A\cap B})=\chi(f_A)\chi(f_B)=1$.
\item[c)] If $A\subset B$ and $A\in\Phi$, then $B\in\Phi$. Indeed, we have $f_A=f_Af_B$, whence $\chi(f_A)=\chi(f_A)\chi(f_B)$, and since $\chi(f_A)=1$, we find $\chi(f_B)=1$.
\item[d)] We have $E\in\Phi$, i.e. $\chi(f_E)=1$, now $f_E=1$, so $\chi(f_E)=1$.
\end{enumerate}
Thus, $\Phi$ is a filter -- we have used only the fact that $A/\mathfrak{m}$ is a ring $\neq0$ (to conclude that $\emptyset\notin\Phi$). Let us prove that if $A/\mathfrak{m}$ is an integral domain, $\Phi$ is an \textit{ultrafilter}, i.e.:
\begin{enumerate}
\item[e)] If $A\notin\Phi$, then $B\defeq E\backslash A\in\Phi$.
\end{enumerate}
Indeed, we have $f_A+f_B=1$, $f_Af_B=0$, whence